\section{Two-scale KAM theorem II: the limit, the tori and the measure}\label{bpP:p3b}

We continue with the notation of Definition~\ref{bpP:p3:seq} and Lemma~\ref{bpP:p3:iterative}; $c:=24^{\tau+1}/6$ and $q:=2\cdot24^{\tau+1}=12c$. Since $\hat c\ge2c$, $\hat c\ge2^{13}n$, the KAM condition gives
\[\hat E_0<\frac1{8192\,n},\qquad c\hat E_0<\frac12.\tag{3.4}\]

\begin{lemma}[Sums]\label{bpP:p3:sums}
For $m\ge0$: $\hat E_k\le\hat E_m^{2^{k-m}}$ for $k\ge m$, $\sum_{k\ge m}\hat E_k\le2\hat E_m$, and $\sum_{k\ge0}c^k\hat E_k\le2\hat E_0$. Also $\hat\rho_k\le\hat\rho_0\,16^{-k}$, $s_k=s\,12^{-k}$.
\end{lemma}
\begin{proof}
\proofstep{Doubling}\label{stepP:p3:sum-doubling}
By Lemma~\ref{bpP:p3:iterative}(iv), $\hat E_{k}<\hat E_{k-1}^2$, so $\hat E_k\le\hat E_m^{2^{k-m}}\le\hat E_m^{\,k-m+1}$ (as $2^{i}\ge i+1$ and $\hat E_m<1$). Hence $\sum_{k\ge m}\hat E_k\le\hat E_m/(1-\hat E_m)\le2\hat E_m$.
\proofstep{Weighted sum}\label{stepP:p3:sum-weighted}
$c^k\hat E_k\le c^k\hat E_0^{k+1}=\hat E_0(c\hat E_0)^k\le\hat E_02^{-k}$ by (3.4); sum $\le2\hat E_0$.
\proofstep{Radii}\label{stepP:p3:sum-radii}
$\hat\rho_k\le\rho_k=\hat\rho_{k-1}/16$ (Definition of $\hat\rho_k$), iterate.
\end{proof}

\subsection{The limit actions (CP10 Step 1)}\label{bpP:p3:limit}

\begin{lemma}[Limit actions]\label{bpP:p3:limit-actions}
Put $\check\imath_0=\mathrm{id}$, $\check\imath_j:=\hat\imath_j\circ\cdots\circ\hat\imath_1:D_0\to D_j$. Then $\check\imath_j$ converges uniformly on $D_0$ to $\check\imath=(\check\imath_1,\check\imath_2)$ and
\begin{enumerate}[label=(\roman*)]
\item $\sup_{D_0}\abs{\check\imath_i-\mathrm{id}_i}\le20\frac{\bar M_i}{\bar M}\hat E_0\hat\rho_0\le20\hat E_0\hat\rho_0$; $\ \sup_{D_0}\abs{\check\imath-\check\imath_j}\le20\hat\rho_j\hat E_j$;
\item $D_*:=\check\imath(D_0)\subset(D_j)_{\rho_j/2}$ for all $j$, in particular $D_*\subset\operatorname{Re}(D_0)_{20\hat E_0\hat\rho_0}$;
\item $\lambda:=L_{D_0}(\check\imath-\mathrm{id})\le2560\,n\hat E_0<\frac13$; so $\check\imath$ is injective, $(1-\lambda)\abs{x-y}\le\abs{\check\imath x-\check\imath y}\le(1+\lambda)\abs{x-y}$;
\item $\omega_j(\check\imath_j(I_0))=\omega_0(I_0)$ for all $I_0\in D_0$, $j\ge0$.
\end{enumerate}
\source{CP10 Step 1, pp.~570--571, (97)--(104); constants recomputed (CP10: $10n$, $2^8n\hat E$).}
\end{lemma}
\begin{proof}
\proofstep{Well-posedness}\label{stepP:p3:la-wp}
$\hat\imath_{j}:D_{j-1}\to D_{j}$ (Lemma~\ref{bpP:p3:iterative}(i)); so $\check\imath_j(D_0)\subset D_j$ and the composition is defined. (iv) follows from $\omega_j\circ\hat\imath_j=\omega_{j-1}$ by induction.
\proofstep{Telescoping}\label{stepP:p3:la-tele}
$\abs{\check\imath_{j,i}(x)-\check\imath_{j-1,i}(x)}=\abs{\hat\imath_{j,i}(y)-y_i}$ with $y=\check\imath_{j-1}(x)\in D_{j-1}$, so $\le10\frac{\bar M_i}{\bar M}\hat E_{j-1}\hat\rho_{j-1}$ (the ratio $\bar M_{i,j}/\bar M_j=\bar M_i/\bar M$ is constant). Since $\hat\rho_{j-1}\le\hat\rho_0$, Lemma~\ref{bpP:p3:sums} gives the Cauchy property and (i): $\sum_{j\ge1}10\hat E_{j-1}\hat\rho_{j-1}\le20\hat E_0\hat\rho_0$; the tail is $\le10\hat\rho_j\sum_{m\ge j}\hat E_m\le20\hat\rho_j\hat E_j$.
\proofstep{Location}\label{stepP:p3:la-loc}
$\check\imath_j(D_0)\subset D_j\subset\R^n$ and $20\hat\rho_j\hat E_j\le\rho_j/2$ ($\hat\rho_j\le\rho_j$, $\hat E_j<1/40$), so $D_*\subset\operatorname{Re}(D_j)_{\rho_j/2}$; for $j=0$ and $D_0\subset D$ this gives the last claim of (ii).
\proofstep{Lipschitz constant}\label{stepP:p3:la-lip}
Put $i_j=\check\imath_j-\mathrm{id}$. From $\check\imath_{j+1}=\hat\imath_{j+1}\circ\check\imath_j$: $i_{j+1}=(\hat\imath_{j+1}-\mathrm{id})\circ\check\imath_j+i_j$, so $1+L(i_{j+1})\le(1+L(\hat\imath_{j+1}-\mathrm{id}))(1+L(i_j))$ and, by Lemma~\ref{bpP:p3:iterative}(i), $1+L(i_j)\le\prod_{k\le j}(1+640n\hat E_{k-1})\le\exp(640n\sum_k\hat E_{k-1})\le\exp(1280n\hat E_0)$. As $x:=1280n\hat E_0<\frac16$, $e^x\le1+2x$: $L(i_j)\le2560n\hat E_0$ uniformly in $j$; the pointwise limit of Lipschitz maps keeps the constant. Injectivity and the two-sided bound follow from $\abs{\check\imath x-\check\imath y}\ge\abs{x-y}-\lambda\abs{x-y}$.
\end{proof}

\subsection{The limit symplectic map (CP10 Step 2)}\label{bpP:p3:limitmap}

\emph{Rescaling.} Let $S(I,\theta):=(I/\hat\rho_0,\ 2\pi\theta/s_0)$ and $\abs{(I,\theta)}_*:=\abs{S(I,\theta)}$ (max norm); for a map $G$ put $\check G=SGS^{-1}$ and $L_*$ for Lipschitz constants in $\abs\cdot_*$. Put $\mathcal D_j:=\operatorname{Re}(D_j)_{\rho_j/2}$. All maps below are real on real points.

\begin{lemma}[Domains and rescaled Lipschitz bounds for $\Psi_k$]\label{bpP:p3:psi-lip}
\begin{enumerate}[label=(\roman*)]
\item $\Psi_k(\mathcal D_k\times\R^n)\subset\mathcal D_{k-1}\times\R^n$ ($k\ge1$); hence $\Phi_j=\Psi_1\circ\dots\circ\Psi_j$ is defined on $\mathcal D_j\times\R^n$ with values in $\mathcal D_0\times\R^n$.
\item $L_*(\check\Psi_k-\mathrm{id};\mathcal D_k\times\R^n)\le192\,n\,c^{k-1}\hat E_{k-1}=:\varsigma_k$, and $\sum_{k\ge1}\varsigma_k\le384n\hat E_0<\frac1{20}$.
\item $\abs{\Psi_k(y)-y}_*\le\frac34\hat E_{k-1}12^{-(k-1)}\le\frac34\hat E_{k-1}$ for $y\in\mathcal D_k\times\R^n$.
\end{enumerate}
\end{lemma}
\begin{proof}
\proofstep{Domains}\label{stepP:p3:pl-dom}
Let $y=(I,\theta)$, $I\in\mathcal D_k$: $\dist(I,D_k)<\rho_k/2=\hat\rho_{k-1}/32$ and $D_k\subset\operatorname{Re}(D_{k-1})_{\hat\rho_{k-1}/16}$. By Lemma~\ref{bpP:p3:iterative}(iii) the real displacement is $\le\frac34\hat E_{k-1}\hat\rho_{k-1}$. Hence $\Psi_k(y)_I$ is within $\hat\rho_{k-1}(\frac1{32}+\frac1{16}+\frac34\hat E_{k-1})<\hat\rho_{k-1}/8$ of $D_{k-1}$; and $\hat\rho_{k-1}/8<\rho_{k-1}/2$ (for $k=1$ because $\hat\rho_0\le\rho$; for $k\ge2$ because $\hat\rho_{k-1}\le\hat\rho_{k-2}/16$).
\proofstep{Bound on displacement}\label{stepP:p3:pl-disp}
Scaled: $\abs{\Delta I}/\hat\rho_0\le\frac34\hat E_{k-1}\hat\rho_{k-1}/\hat\rho_0\le\frac34\hat E_{k-1}16^{-(k-1)}$ (Lemma~\ref{bpP:p3:sums}; $\hat M\le M$) and $2\pi\abs{\Delta\theta}/s_0\le\frac34\hat E_{k-1}s_{k-1}/s_0=\frac34\hat E_{k-1}12^{-(k-1)}$. Thus $A_k:=\sup_{W_k}\abs{\check\Psi_k-\mathrm{id}}\le\frac34\hat E_{k-1}12^{-(k-1)}$; this is (iii).
\proofstep{Lipschitz}\label{stepP:p3:pl-lip}
Apply Lemma~\ref{bpP:p3:lip-sup} in the scaled coordinates, $N=2n$, $S=S(\mathcal D_k\times\R^n)$, radius $\varepsilon_k=\min\{\rho_k/(2\hat\rho_0),\ s_k/s_0\}$: the scaled complex $\varepsilon_k$-neighbourhood of a real point of $S$ lies in $(D_k)_{\rho_k}\times\T_{s_k}\subset W_k$ ($\varepsilon_k\hat\rho_0\le\rho_k/2$ and $\varepsilon_ks_0\le s_k\le\bar s+s_k$), where the bound $A_k$ of (iii) holds (Lemma~\ref{bpP:p3:iterative}(iii) holds on the complex $W_k$). Now $\rho_k/(2\hat\rho_0)=\hat\rho_{k-1}/(32\hat\rho_0)\ge q^{-(k-1)}/32$ (Lemma~\ref{bpP:p3:iterative}(v)) and $s_k/s_0=12^{-k}\ge q^{-(k-1)}/32$, so $\varepsilon_k\ge q^{-(k-1)}/32$ and
\[L_*\le\frac{4\cdot2n\,A_k}{\varepsilon_k}\le8n\cdot\tfrac34\hat E_{k-1}12^{-(k-1)}\cdot32q^{k-1}=192n\,(q/12)^{k-1}\hat E_{k-1}=192n\,c^{k-1}\hat E_{k-1}.\]
\proofstep{Sum}\label{stepP:p3:pl-sum}
Lemma~\ref{bpP:p3:sums}: $\sum_k\varsigma_k=192n\sum_{m\ge0}c^m\hat E_m\le384n\hat E_0<384/8192<\frac1{20}$.
\end{proof}

\begin{lemma}[Compactness of periodic holomorphic functions on a strip]\label{bpP:p3:fourier-compact}
Let $F_j:\{\vartheta\in\C^n:\abs{\operatorname{Im}\vartheta}<w\}\to\C$ be holomorphic, $\Z^n$-periodic, $\abs{F_j}\le B$, and suppose $F_j$ converges uniformly on $\R^n$ to $F$. Then for \emph{every} $0<v<w$, $F_j$ converges uniformly, together with all derivatives, on $\{\abs{\operatorname{Im}\vartheta}\le v\}$; the limit is a holomorphic periodic function on the whole open strip extending $F$, with $\abs F\le B$ there.
\end{lemma}
\begin{proof}
\proofstep{Coefficient decay}\label{stepP:p3:fc-decay}
For $0<w'<w$ and $k\ne0$ shift the integration contour to $\operatorname{Im}\vartheta=-w'\operatorname{sign}(k)$ (componentwise; Cauchy's theorem for the periodic integrand, \Kref{bp:p1:ext-rectangle}): $\abs{\hat F_{j,k}}\le Be^{-2\pi w'\abs k_1}$, uniformly in $j$.
\proofstep{Convergence of coefficients}\label{stepP:p3:fc-coef}
Uniform convergence on $\R^n$ gives $\hat F_{j,k}\to\hat F_k$ for every $k$ (\ref{bpP:p3:ext-fourier}), and $\abs{\hat F_k}\le Be^{-2\pi w'\abs k_1}$ for every $w'<w$.
\proofstep{Tail and finite part on $\abs{\operatorname{Im}\vartheta}\le v$}\label{stepP:p3:fc-tail}
Fix $0<v<w$ and put $w'=(v+w)/2>v$. There $\abs{e^{2\pi ik\cdot\vartheta}}\le e^{2\pi v\abs k_1}$, so for every $m\ge0$
$\sum_{\abs k_1>N}\abs{(2\pi k)^{m}\hat F_{j,k}e^{2\pi ik\cdot\vartheta}}\le B\sum_{\abs k_1>N}(2\pi\abs k_1)^me^{-2\pi(w'-v)\abs k_1}=:\varepsilon_{N,m}\to0$ ($N\to\infty$), uniformly in $j$, and the same for $F$; the finitely many modes $\abs k_1\le N$ converge uniformly. Hence $F_j\to F^\sharp:=\sum_k\hat F_ke^{2\pi ik\cdot\vartheta}$ with all derivatives, uniformly on $\abs{\operatorname{Im}\vartheta}\le v$. The limits for different $v$ coincide (same coefficients), so $F^\sharp$ is defined and holomorphic on the whole open strip (\ref{bpP:p3:ext-weierstrass}), equals $F$ on $\R^n$ and, as a pointwise limit of the $F_j$, satisfies $\abs{F^\sharp}\le B$.
\end{proof}

\begin{lemma}[The limit map $\Phi$]\label{bpP:p3:limit-phi}
On $W_*:=D_*\times\R^n$ the maps $\Phi_j$ converge uniformly to a map $\Phi$ with $\Phi-\mathrm{id}$ $\Z^n$-periodic in $\theta$, real analytic in $\theta$, and
\begin{enumerate}[label=(\roman*)]
\item $\sup_{W_*}\abs{\Pi_{I_1}\Phi-\mathrm{id}}\le\frac32\frac{\hat M}M\hat E_0\hat\rho_0$, $\sup\abs{\Pi_{I_2}\Phi-\mathrm{id}}\le\frac32\hat E_0\hat\rho_0$, $\sup2\pi\abs{\Pi_\theta\Phi-\mathrm{id}}\le\frac32\hat E_0s_0$, and the same bounds on complex $\theta$ with $\abs{\operatorname{Im}\theta}<\bar s/2\pi$;
\item $L_*(\check\Phi-\mathrm{id};W_*)\le768\,n\hat E_0<\frac1{10}$; hence $\Phi$ is injective on $W_*$ and $\abs{x-y}_*\cdot\frac9{10}\le\abs{\Phi x-\Phi y}_*\le\frac{11}{10}\abs{x-y}_*$;
\item $\Phi(I,\theta+k)=\Phi(I,\theta)+(0,k)$ for $k\in\Z^n$.
\end{enumerate}
\source{CP10 Step 2, pp.~571--572, (105)--(108).}
\end{lemma}
\begin{proof}
\proofstep{Telescoping of the displacements}\label{stepP:p3:lp-tele}
Write $(I^{(k)},\theta^{(k)})=\Psi_{k+1}\circ\dots\circ\Psi_j(x)$ for $x\in W_*\subset W_j$ (Lemma~\ref{bpP:p3:limit-actions}(ii)); then $\Pi_I\Phi_j(x)-\Pi_Ix=\sum_{k=1}^{j}(I^{(k-1)}-I^{(k)})$ and every summand is bounded by Lemma~\ref{bpP:p3:iterative}(iii). Summing with Lemma~\ref{bpP:p3:sums} ($\hat\rho_{k-1}\le\hat\rho_0$, $s_{k-1}\le s_0$, $\sum\hat E_{k-1}\le2\hat E_0$) gives (i), uniformly in $j$ and also for complex $\theta$ ($\Psi_k$ is defined on the complex strip).
\proofstep{Cauchy sequence}\label{stepP:p3:lp-cauchy}
For $x\in W_*$: $\Phi_j(x)-\Phi_{j-1}(x)=\Phi_{j-1}(\Psi_jx)-\Phi_{j-1}(x)$, with $x\in\mathcal D_j\times\R^n$ and $\Psi_jx\in\mathcal D_{j-1}\times\R^n$ (Lemma~\ref{bpP:p3:psi-lip}(i)). By the composition rule $L_*(\check\Phi_{j-1}-\mathrm{id})\le\prod_{k<j}(1+\varsigma_k)-1\le e^{\sum\varsigma_k}-1\le2\sum\varsigma_k\le\frac1{10}$ (using $e^x\le1+2x$, $x\le1$; the composition formula is that of Lemma~\ref{bpP:p3:limit-actions}, proof step~\ref{stepP:p3:la-lip}, valid because each $\Psi_k$ maps its domain into the next one). Therefore $\abs{\Phi_j(x)-\Phi_{j-1}(x)}_*\le\frac{11}{10}\abs{\Psi_jx-x}_*\le\frac{11}{10}\cdot\frac34\hat E_{j-1}$, which is summable: uniform convergence on $W_*$. The limit satisfies (ii) with $L_*\le2\sum\varsigma_k\le768n\hat E_0$ (the constant passes to the pointwise limit).
\proofstep{Periodicity and injectivity}\label{stepP:p3:lp-per}
Each $\Psi_k-\mathrm{id}$ is $\Z^n$-periodic in $\theta$ (Lemma~\ref{bpP:p3:iterative}), so (iii) holds for $\Phi_j$ and in the limit. From (ii), $\abs{\Phi x-\Phi y}_*\ge\frac9{10}\abs{x-y}_*$.
\proofstep{Analyticity in $\theta$}\label{stepP:p3:lp-an}
Fix $I_*\in D_*$; $\vartheta\mapsto\Phi_j(I_*,\vartheta)-(0,\vartheta)$ is holomorphic on $\{\abs{\operatorname{Im}\vartheta}<\bar s/2\pi\}$ (as $I_*\in(D_j)_{\rho_j}$), $\Z^n$-periodic, bounded uniformly in $j$ by (i), and converges uniformly on $\R^n$. By Lemma~\ref{bpP:p3:fourier-compact} (componentwise, $w=\bar s/2\pi$) it converges, with all derivatives, uniformly on $\{\abs{\operatorname{Im}\vartheta}\le v\}$ for every $v<\bar s/2\pi$, to a function holomorphic on the \emph{whole} strip $\{\abs{\operatorname{Im}\vartheta}<\bar s/2\pi\}$ whose real restriction is $\Phi(I_*,\cdot)-(0,\mathrm{id})$; the bounds (i), being pointwise, pass to the limit. This is claim (a) of Proposition~\ref{bpP:p3:two-scale-kam}.
\end{proof}

\subsection{The tori (CP10 Steps 2--3 and the estimates (35))}\label{bpP:p3:tori}

Define, for $\omega\in\Omega_*=\omega_0(D)\cap\mathcal D_{\gamma_1,\gamma_2,\tau}$, $I_0(\omega):=\omega_0^{-1}(\omega)\in D_0$ (injectivity of $\omega_0$; $\omega\in\omega_0(D)\cap\mathcal D$ means $I_0\in D^{\rm dio}=D_0$), $I_*(\omega):=\check\imath(I_0(\omega))\in D_*$ and
\[\varphi_\omega(\vartheta):=\Phi(I_*(\omega),\vartheta)=(v(\vartheta;\omega),\ \vartheta+u(\vartheta;\omega)).\]

\begin{lemma}[Estimates and location]\label{bpP:p3:estimates}
For $\omega\in\Omega_*$: $\abs{u}\le\frac3{4\pi}\hat E_0s$ on $\{\abs{\operatorname{Im}\vartheta}<\bar s/2\pi\}$; $\abs{v_1-I_{0,1}}\le20(\frac{\bar M_1}{\bar M}+\frac{\hat M}M)\hat E_0\hat\rho_0$, $\abs{v_2-I_{0,2}}\le20(\frac{\bar M_2}{\bar M}+1)\hat E_0\hat\rho_0$ on $\R^n$, and $v\in\operatorname{Re}(D_{r})$ with $r=32\hat E_0\hat\rho_0$, so $\varphi_\omega(\R^n)\subset\operatorname{Re}(D_r)\times\R^n$.
This is Proposition~\ref{bpP:p3:two-scale-kam}(a), (d), and the location claim.
\source{CP10 (35), (105)--(107), Step 2 end, p.~572. \emph{Misprint:} CP10 (35) bounds $\abs{v_2-I_2^0}$ by $10n(\bar M_2/\bar M+\hat M/M)\hat E\hat\rho$ although (107) only gives $2\hat E\hat\rho$ for $I_2$ without the factor $\hat M/M$.}
\end{lemma}
\begin{proof}
\proofstep{Sup bounds}\label{stepP:p3:es-sup}
$\abs{v_i-I_{0,i}}\le\abs{\Pi_{I_i}\Phi(I_*,\vartheta)-I_{*,i}}+\abs{I_{*,i}-I_{0,i}}$, the two terms bounded by Lemma~\ref{bpP:p3:limit-phi}(i) and Lemma~\ref{bpP:p3:limit-actions}(i): $\frac32\frac{\hat M}M\hat E\hat\rho+20\frac{\bar M_1}{\bar M}\hat E\hat\rho$ for $i=1$ (and $\frac32+20\frac{\bar M_2}{\bar M}$ for $i=2$), both $\le$ the stated constants. $\abs u\le\frac32\hat Es/2\pi$ is the angle bound.
\proofstep{Location}\label{stepP:p3:es-loc}
$I_*\in\operatorname{Re}(D_0)_{20\hat E\hat\rho}$ and $\abs{v-I_*}\le\frac32\hat E\hat\rho$, so $\dist(v,D)<22\hat E\hat\rho<r$.
\end{proof}

\begin{lemma}[Immersion and Lagrangian property]\label{bpP:p3:lagr}
For $\omega\in\Omega_*$, $\norm{D_\vartheta u}\le\frac34n\hat E_0$ (so $D_\vartheta\varphi_\omega=\binom{D_\vartheta v}{1+D_\vartheta u}$ has injective lower block $1+D_\vartheta u$, since $\frac34n\hat E_0<\frac12$), and $\varphi_\omega^*(\sum_idJ_i\wedge d\theta_i)=0$.
\end{lemma}
\begin{proof}
\proofstep{Cauchy in the strip}\label{stepP:p3:lg-cauchy}
$u$ is holomorphic on $\abs{\operatorname{Im}\vartheta}<\bar s/2\pi$ with $\abs u\le3\hat Es/4\pi$. At real $\vartheta$ the polydisc of radius $\bar s/4\pi$ lies in the strip, so $\abs{\partial_au_i}\le\frac{3\hat Es}{4\pi}\cdot\frac{4\pi}{\bar s}=3\hat E\frac s{\bar s}\le\frac34\hat E$ ($4s\le\bar s$), and the row sum is $\le\frac34n\hat E$.
\proofstep{Each $\Phi_j$-torus is Lagrangian}\label{stepP:p3:lg-j}
$\Phi_j$ is symplectic and $\vartheta\mapsto(I_*,\vartheta)$ has pullback $\sum dI_i\wedge d\theta_i=0$ (the $I$-coordinates are constant), so $\varphi^{(j)}:=\Phi_j(I_*,\cdot)$ has $\Omega_j(a,b):=\sum_i(\partial_av^{(j)}_i\partial_b\Theta^{(j)}_i-\partial_bv_i^{(j)}\partial_a\Theta_i^{(j)})=0$.
\proofstep{Limit}\label{stepP:p3:lg-lim}
By Lemma~\ref{bpP:p3:limit-phi} (Step \ref{stepP:p3:lp-an}) the first derivatives of $\varphi^{(j)}$ converge to those of $\varphi_\omega$ (uniformly on $\R^n$, Lemma~\ref{bpP:p3:fourier-compact}); $\Omega_j$ is bilinear in them, so $\Omega=\lim\Omega_j=0$.
\end{proof}

\begin{lemma}[Invariance]\label{bpP:p3:invariance}
For $\omega\in\Omega_*$, $\vartheta\in\R^n$, $t\in\R$: $\frac{d}{dt}\varphi_\omega(\vartheta+t\omega)=X_H(\varphi_\omega(\vartheta+t\omega))$, $X_H=(-\partial_\theta H,\partial_JH)$.
\source{CP10 Step 3, pp.~572--573.}
\end{lemma}
\begin{proof}
Let $I_0=I_0(\omega)$, $I_*=\check\imath(I_0)$, $\theta_t=\vartheta+t\omega$, $\gamma(t)=\Phi(I_*,\theta_t)$, $\gamma_j(t)=\Phi_j(I_*,\theta_t)$.
\proofstep{Symplectic conjugation}\label{stepP:p3:inv-conj}
Since $\Phi_j$ is a symplectic diffeomorphism onto its image and $H_j=H\circ\Phi_j$: $X_H\circ\Phi_j=D\Phi_j\,X_{H_j}$ (\Nref{bpN:p5:symplectic}). With $X_{H_j}(I_*,\theta)=(-\partial_\theta f_j,\ \omega_j(I_*)+\partial_If_j)$ we get
$\gamma_j'(t)-X_H(\gamma_j(t))=D\Phi_j(I_*,\theta_t)\bigl(\partial_\theta f_j,\ \omega-\omega_j(I_*)-\partial_If_j\bigr)$ (with $\gamma_j'=D\Phi_j(0,\omega)$).
\proofstep{$D\Phi_j$ bounded}\label{stepP:p3:inv-dphi}
$\mathcal D_j\times\R^n$ is open and $L_*(\check\Phi_j)\le\frac{11}{10}$ (Lemma~\ref{bpP:p3:limit-phi}(ii) argument), so $\abs{S\,D\Phi_jS^{-1}}\le\frac{11}{10}$: $\norm{D\Phi_j}\le\kappa:=\frac{11}{10}\max\{1,\tfrac{2\pi\hat\rho_0}{s_0},\tfrac{s_0}{2\pi\hat\rho_0}\}$, independent of $j$.
\proofstep{The three error terms}\label{stepP:p3:inv-err}
(1) $\omega=\omega_j(\check\imath_j(I_0))$ (Lemma~\ref{bpP:p3:limit-actions}(iv)); $\abs{I_*-\check\imath_j(I_0)}\le20\hat\rho_j\hat E_j$ and the segment lies in $(D_j)_{\rho_j}$ where $\norm{U_j}\le M_j$, so $\abs{\omega-\omega_j(I_*)}\le20M_j\hat\rho_j\hat E_j\le7\gamma_1\hat E_j$, using $M_j\hat\rho_j\le\gamma_1/(3K_j^{\tau+1})\le\gamma_1/3$.
(2) $I_*\in\operatorname{Re}(D_j)_{\rho_j/2}$, so the polydisc of radius $\rho_j/2$ about $I_*$ is in $(D_j)_{\rho_j}$ and \Kref{bp:p1:cauchy-open} gives $\abs{\partial_If_j}\le2E_j/\rho_j\le2E_j/\hat\rho_j=2\hat E_j\hat\rho_j/L_j\le2\hat E_jM_j\hat\rho_j\le\gamma_1\hat E_j$ ($L_j\ge M_j^{-1}$).
(3) For real $\theta$ the polydisc of radius $\bar s/4\pi$ in $\theta$ lies in $\T_{\bar s/2}$, where $\abs{f_j}\le E_j$ (as $\bar s/2\le\bar s+s_j$), so $\abs{\partial_\theta f_j}\le4\pi E_j/\bar s$ (\Kref{bp:p1:cauchy-open}); and $E_j=\hat E_j\hat\rho_j^2/L_j\le\hat E_j\hat\rho_jM_j\hat\rho_j\le\hat E_j\gamma_1\hat\rho_0$.
All three tend to $0$, uniformly in $(\vartheta,t)$, as $\hat E_j\le\hat E_0^{2^j}$. Hence $\sup_t\abs{\gamma_j'-X_H(\gamma_j)}\to0$.
\proofstep{Passing to the limit}\label{stepP:p3:inv-lim}
$\delta_j:=\sup_t\abs{\gamma_j(t)-\gamma(t)}\to0$ (Lemma~\ref{bpP:p3:limit-phi}; unscaled distances are comparable to scaled ones with fixed constants). Only nearby pairs are needed (review P3-08). The action of $\gamma(t)$ lies in $\operatorname{Re}(D_r)\subset\operatorname{Re}(D_{\rho/4})$ ($r=32\hat E\hat\rho\le\rho/4$ since $\hat E<2^{-13}$, $\hat\rho\le\rho$) and its angle is real; for $j$ so large that $\delta_j<\min\{\rho/4,\bar s/4\pi\}$, the real segment $[\gamma_j(t),\gamma(t)]$ lies in the convex polydisc of radius $\delta_j$ about $\gamma(t)$, hence in $D_{\rho/2}\times\T_{\bar s/2}$. There $\norm{DX_H}\le B_0$, where $DX_H$ consists of $U+\partial_J^2f$, $\partial_\theta\partial_Jf$, $\partial_\theta^2f$ and $B_0=M+C(E,\rho,\bar s)$ by \Kref{bp:p1:cauchy-open} (applied to $\sup\abs f\le E$ on $D_\rho\times\T_{\bar s}$). By the mean value inequality, $\abs{X_H(\gamma_j(t))-X_H(\gamma(t))}\le B_0\delta_j\to0$ uniformly in $t$. A uniform limit of $C^1$ functions $\gamma_j\to\gamma$ with $\gamma_j'\to X_H(\gamma)$ uniformly is $C^1$ with $\gamma'=X_H(\gamma)$. Since $\gamma(t)=\varphi_\omega(\vartheta+t\omega)$, this is the claim. (Conjugation identity of Step~\ref{stepP:p3:inv-conj}: \Nref{bpN:p5:symplectic}, \Nref{bpN:p5:canonical}.)
\end{proof}

\subsection{Measure estimate (CP10 Step 4)}\label{bpP:p3:measure}

Let $K_\Omega=\bigcup_{\omega\in\Omega_*}\varphi_\omega(\R^n)$; by definition $K_\Omega=\Phi(D_*\times\R^n)$ ($\check\imath$ maps $D_0=D^{\rm dio}$ onto $D_*$ and $\omega\mapsto I_0(\omega)$ is a bijection $\Omega_*\to D_0$). Here and below $N_\rho(X)$ is the \emph{closed} real max-norm $\rho$-neighbourhood $\{x:\dist_\infty(x,X)\le\rho\}$; every radius used is $<r$, so $N_\rho(X)\subset\operatorname{Re}(X_r)$. Put $a:=\meas(D\setminus D_0)$, $b:=\meas(\operatorname{Re}(D_r)\setminus D)$, $\lambda\le\frac13$ as in Lemma~\ref{bpP:p3:limit-actions}, $\eta:=\sup\abs{\check\imath-\mathrm{id}}\le20\hat E_0\hat\rho_0$.

\begin{lemma}[Bi-Lipschitz extension and covering]\label{bpP:p3:cover}
Let $X\subset\R^N$ be arbitrary, $g:X\to\R^N$ with $L(g)\le\lambda<1$ in the max norm and $\sup\abs g\le\eta$. Then there is a bijection $F=\mathrm{id}+g_e:\R^N\to\R^N$, $F|_X=\mathrm{id}+g$, with $(1-\lambda)\abs{x-y}\le\abs{F(x)-F(y)}\le(1+\lambda)\abs{x-y}$, and for $y\in X$: $\dist_\infty(F^{-1}(y),X)\le\abs{F^{-1}(y)-y}\le\eta/(1-\lambda)$ (a distance inequality, valid also for $\eta=0$).
\end{lemma}
\begin{proof}
\proofstep{Extension}\label{stepP:p3:cv-ext}
$g_e$: McShane \Kref{bp:p1:mcshane}, componentwise, with the same $\lambda$ for the max norm.
\proofstep{Bijectivity}\label{stepP:p3:cv-bij}
$\abs{F(x)-F(y)}\ge\abs{x-y}-\lambda\abs{x-y}$ gives injectivity and the lower bound; for $y\in\R^N$ the map $x\mapsto y-g_e(x)$ is a $\lambda$-contraction (\ref{bpP:p3:ext-banach}), with fixed point $x=F^{-1}(y)$.
\proofstep{Location of preimages}\label{stepP:p3:cv-loc}
For $y\in X$ and $x=F^{-1}(y)$: $\abs{x-y}=\abs{g_e(x)}\le\abs{g_e(y)}+\lambda\abs{x-y}\le\eta+\lambda\abs{x-y}$ so $\abs{x-y}\le\eta/(1-\lambda)$.
\end{proof}

\begin{lemma}[Measure of the complement of $D_*$]\label{bpP:p3:meas-actions}
$\meas(D_0\setminus D_*)\le(4/3)^n(a+b)$ and hence $\meas(D\setminus D_*)\le a+(4/3)^n(a+b)$.
\end{lemma}
\begin{proof}
\proofstep{Extension and preimage}\label{stepP:p3:ma-ext}
Apply Lemma~\ref{bpP:p3:cover} with $X=D_0$, $g=\check\imath-\mathrm{id}$, $N=n$: $F=\check\imath_e$. Let $\rho_1:=\eta/(1-\lambda)\le30\hat E_0\hat\rho_0<r$. If $y\in D_0\setminus D_*$, then $x=F^{-1}(y)\in N_{\rho_1}(D_0)$ and $x\notin D_0$ (else $y=F(x)=\check\imath(x)\in D_*$). So $D_0\setminus D_*\subset F\bigl(N_{\rho_1}(D_0)\setminus D_0\bigr)$.
\proofstep{Measure}\label{stepP:p3:ma-meas}
By \Kref{bp:p1:lip-image} ($F$ is $(1+\lambda)$-Lipschitz), $\meas(D_0\setminus D_*)\le(1+\lambda)^n\meas(N_{\rho_1}(D_0)\setminus D_0)$, $(1+\lambda)\le\frac43$. As $D_0\subset D$: $N_{\rho_1}(D_0)\setminus D_0\subset(N_{\rho_1}(D)\setminus D)\cup(D\setminus D_0)$, of measure $\le b+a$ ($\rho_1<r$). The second claim: $D\setminus D_*\subset(D\setminus D_0)\cup(D_0\setminus D_*)$.
\end{proof}

\begin{lemma}[Measure of the complement of $K_\Omega$ over $D$]\label{bpP:p3:meas-torus}
$\meas\bigl((D\times\T^n)\setminus K_\Omega\bigr)\le\bigl[1+(4/3)^n+2^n\bigl(1+(4/3)^n\bigr)\bigr](a+b)$.
\end{lemma}
\begin{proof}
Identify $\T^n$-measure with the unit cell; $K_\Omega$ is $\Z^n$-invariant in the angle.
\proofstep{Splitting}\label{stepP:p3:mt-split}
$(D\times\T^n)\setminus K_\Omega\subset\bigl((D\setminus D_*)\times\T^n\bigr)\cup\bigl(((D_*\cap D)\times\T^n)\setminus K_\Omega\bigr)$. The first set has measure $\le a+(4/3)^n(a+b)$ by Lemma~\ref{bpP:p3:meas-actions}.
\proofstep{Rescaled extension}\label{stepP:p3:mt-ext}
Work in scaled coordinates $x\mapsto Sx$; $S$ is linear, multiplies the Lebesgue measure of every subset of $\R^n_I\times\R^n_\theta$ by the constant $\hat\rho_0^{-n}(2\pi/s_0)^n$ and maps products to products, so all comparisons below between measures of subsets of $\R^n\times\R^n$ are unaffected. By Lemma~\ref{bpP:p3:limit-phi}, $g=\check\Phi-\mathrm{id}$ on $X=S(D_*\times\R^n)\subset\R^{2n}$ has $L_*(g)\le\lambda':=768n\hat E_0<\frac1{10}$ and $\sup\abs g\le\eta':=\frac32\hat E_0$ (Lemma~\ref{bpP:p3:limit-phi}(i), scaled). Lemma~\ref{bpP:p3:cover} ($N=2n$) gives the bijection $\check\Phi_e=\mathrm{id}+g_e$ of $\R^{2n}$ extending $\check\Phi$ on $X$, $(1+\lambda')$-Lipschitz, with $\abs{\check\Phi_e^{-1}(y)-y}_*\le\eta'/(1-\lambda')\le2\hat E_0$ for $y\in X$.
\proofstep{Preimage structure}\label{stepP:p3:mt-pre}
Let $y=(I_y,\vartheta_y)$ with $I_y\in D_*\cap D$ and $\vartheta_y\in[0,1)^n$ be a point of the second set, i.e.\ $y\notin K_\Omega$ (mod $\Z^n$ in the angle), and let $x=(I_x,\vartheta_x):=\Phi_e^{-1}(y)$ (unscaled notation; $y\in D_*\times\R^n=S^{-1}X$). Then $\abs{I_x-I_y}\le\rho_2:=2\hat E_0\hat\rho_0$ and $\abs{\vartheta_x-\vartheta_y}\le\hat E_0s_0/\pi\le\hat E_0/(4\pi)=:\epsilon$ (as $s_0\le\frac14$). If $I_x\in D_*$, then $x\in D_*\times\R^n$, $\Phi_e(x)=\Phi(x)=y$ and $y\in K_\Omega$, a contradiction. Hence $I_x\in N_{\rho_2}(D_*)\setminus D_*$ and $\vartheta_x\in[-\epsilon,1+\epsilon)^n$.
\proofstep{Measure}\label{stepP:p3:mt-meas}
Thus the second set is contained in $\Phi_e\bigl((N_{\rho_2}(D_*)\setminus D_*)\times[-\epsilon,1+\epsilon)^n\bigr)$. $\Phi_e$ is $(1+\lambda')$-Lipschitz for $\abs\cdot_*$ on $\R^{2n}$, so by \Kref{bp:p1:lip-image} in scaled coordinates (the scaling constants cancel on both sides) its measure is $\le(1+\lambda')^{2n}(1+2\epsilon)^n\meas(N_{\rho_2}(D_*)\setminus D_*)\le2^n\meas(N_{\rho_2}(D_*)\setminus D_*)$, because $(1.1)^{2n}(1+\hat E_0/2\pi)^n\le(1.21\cdot1.01)^n<2^n$.
\proofstep{The shell}\label{stepP:p3:mt-shell}
$D_*\subset N_\eta(D_0)$, so $N_{\rho_2}(D_*)\subset N_{\rho_2+\eta}(D_0)\subset N_{\rho_2+\eta}(D)\subset\operatorname{Re}(D_r)$ ($\rho_2+\eta\le22\hat E_0\hat\rho_0<r$), and $N_{\rho_2}(D_*)\setminus D_*\subset(N_{\rho_2+\eta}(D)\setminus D)\cup(D_0\setminus D_*)\cup(D\setminus D_0)$: if $x\in N_{\rho_2}(D_*)\setminus D_*$ then either $x\notin D$ (first set), or $x\in D\setminus D_0$, or $x\in D_0\setminus D_*$. Its measure is $\le b+a+(4/3)^n(a+b)=(1+(4/3)^n)(a+b)$ by Lemma~\ref{bpP:p3:meas-actions}.
\proofstep{Sum}\label{stepP:p3:mt-sum}
Total $\le[a+(4/3)^n(a+b)]+2^n(1+(4/3)^n)(a+b)\le[1+(4/3)^n+2^n(1+(4/3)^n)](a+b)$.
\end{proof}

\subsection{Proof of Proposition~\ref{bpP:p3:two-scale-kam}}\label{bpP:p3:proof}

\begin{proof}[Proof of Proposition~\ref{bpP:p3:two-scale-kam}]
\proofstep{Construction}\label{stepP:p3:pf-constr}
The sequence of Lemma~\ref{bpP:p3:iterative} exists (hypotheses (S1)--(S3) at $j=0$ from the proposition's hypotheses and the definition of $D_0$); Lemma~\ref{bpP:p3:limit-actions} defines $\check\imath$, Lemma~\ref{bpP:p3:limit-phi} defines $\Phi$, and $\varphi_\omega$ is defined in \S\ref{bpP:p3:tori}.
\proofstep{(a), (d)}\label{stepP:p3:pf-ad}
Lemmas~\ref{bpP:p3:limit-phi}, \ref{bpP:p3:estimates}.
\proofstep{(b), (c)}\label{stepP:p3:pf-bc}
Lemmas~\ref{bpP:p3:invariance}, \ref{bpP:p3:lagr}.
\proofstep{Injectivity and Lipschitz structure}\label{stepP:p3:pf-inj}
If $\varphi_\omega(\vartheta)-\varphi_{\omega'}(\vartheta')=(0,k)$, then $\Phi(I_*,\vartheta)=\Phi(I_*',\vartheta'+k)$ by periodicity (Lemma~\ref{bpP:p3:limit-phi}(iii)), so $(I_*,\vartheta)=(I_*',\vartheta'+k)$ by (ii); $I_*=I_*'$ gives $I_0=I_0'$ ($\check\imath$ injective), so $\omega=\omega_0(I_0)=\omega'$. The factorization and the constants: Lemma~\ref{bpP:p3:limit-actions}(iii) and Lemma~\ref{bpP:p3:limit-phi}(ii). For $\omega_0^{-1}$: if $\omega,\omega'\in\omega_0(D)$ with $\abs{\omega-\omega'}<\rho/\bar M$, then Lemma~\ref{bpP:p3:ball} (with $F=\omega_0$ on $D_\rho$, $\delta=\rho$, $x_0=\omega_0^{-1}(\omega)$) gives $\abs{\omega_0^{-1}\omega-\omega_0^{-1}\omega'}\le\bar M\abs{\omega-\omega'}$; for far pairs $\abs{\omega_0^{-1}\omega-\omega_0^{-1}\omega'}\le\operatorname{diam}D\le\operatorname{diam}D\cdot\frac{\bar M}{\rho}\abs{\omega-\omega'}$.
\proofstep{Measure}\label{stepP:p3:pf-meas}
$K_\Omega\subset\operatorname{Re}(D_r)\times\T^n$ (Lemma~\ref{bpP:p3:estimates}), so $(\operatorname{Re}(D_r)\times\T^n)\setminus K_\Omega\subset((\operatorname{Re}(D_r)\setminus D)\times\T^n)\cup((D\times\T^n)\setminus K_\Omega)$ (an inclusion; the shell may meet $K_\Omega$, which only helps), and Lemma~\ref{bpP:p3:meas-torus} gives (3.1) with $c_n=1+[1+(4/3)^n+2^n(1+(4/3)^n)]=1+(1+2^n)(1+(4/3)^n)$ ($a=\meas(D\setminus D^{\rm dio})$, $b=\meas(\operatorname{Re}(D_r)\setminus D)$, $\meas\T^n=1$).
\end{proof}

\subsection{What Part 5 consumes}\label{bpP:p3:interface}

Part~5 (\Pref{bpP:p5:kam-application}) applies Proposition~\ref{bpP:p3:two-scale-kam} to the Hamiltonian $\breve H$ of CP10 Step~6 (p.~560) after Part~2's averaging (\Pref{bpP:p2:averaged-remainder}). \emph{Three radii must not be confused:} $\epsilon$ is the oscillator radius of CP10 (the quantity renamed $r$ in the frozen cylinder $\mathcal C(V,r)$); $\hat\rho$ is the action radius of Definition~\ref{bpP:p3:constants}; $r=32\hat E\hat\rho$ is the enlargement radius of this proposition. The identification (CP10 p.~560, review P3-05):
$n_1=a$, $n_2=b$, $h=H_0(J_1)+\mu\breve N(J)$, $f=\mu\,\epsilon^5\breve P$ (\emph{not} $\mu^5\breve P$), $\rho=\breve\rho$, $s=\breve s/5$, $\bar s=4\breve s/5$ (in coefficient-radius widths of \S\ref{bpP:p3:conv}), $D=D_{\rm CP}$ of CP10 (48).
\begin{enumerate}[label=(\arabic*)]
\item \emph{Before $\epsilon,\mu$} (depend only on the fixed family): $n,\tau$ (hence $\hat c,c_n$), and a constant $C\ge1$ such that the true bounds satisfy $\sup\norm U\le C$, $\sup\norm{\hat U}\le C\mu$, $\sup\norm{U^{-1}}\le C\mu^{-1}$, $\sup\norm{T_1}\le C$, $\sup\norm{T_2}\le C\mu^{-1}$, $\norm f\le C\mu\epsilon^5$ for $0<\mu\le1$, $\epsilon\le1$ (CP10 p.~560; the N-norm of $\mu\epsilon^5\breve P$ is $\le C\mu\epsilon^5$).
\item \emph{Choice of the numerical bounds (review R2-01).} Since $M,\hat M,\bar M,\bar M_i,E$ need only be \emph{upper} bounds, take them \emph{equal} to these values: $M=C$, $\hat M=C\mu$, $\bar M=\bar M_2=C\mu^{-1}$, $\bar M_1=C$, $E=C\mu\epsilon^5$ (the normalizations $\hat M\le M$, $\bar M_i\le\bar M$ hold for $\mu\le1$, $C\ge1$). Then, exactly, $L=\max\{C\mu^{-1},C^{-1},(C\mu)^{-1}\}=C\mu^{-1}$ (because $C\ge C^{-1}$ gives $C\mu^{-1}\ge(C\mu)^{-1}$, and $C\mu^{-1}\ge C^{-1}$ as $\mu\le1$), $EL=C^2\epsilon^5$ and
$x=EM^2L/\gamma_1^2=C^4\epsilon^5/\gamma_1^2$, so $K=\frac6s\log_+\bigl(\gamma_1^2/(C^4\epsilon^5)\bigr)$ \emph{exactly} (no inequality is inferred from an upper bound on $x$: $\log_+(1/x)$ is nonincreasing in $x$, so an upper bound on $x$ would only give a lower bound on $K$). Whenever $\gamma_1^2/(C^4\epsilon^5)\ge e$ this is $K=\frac6s\log(\gamma_1^2/(C^4\epsilon^5))\le C'\log(\epsilon^5/\gamma_1^2)^{-1}$ (CP10 p.~560), with $C'$ depending on $s,C$ only. The smaller true $E$ can only be used if one re-chooses the numerical $E$ \emph{before} computing $K$; Part~5 should keep the choice above. Then $\hat\rho=\min\{\gamma_1/(3CK^{\tau+1}),\gamma_2/(3C\mu K^{\tau+1}),\rho\}$, $\hat E=C^2\epsilon^5/\hat\rho^2$, and the KAM condition $\hat c\hat E<1$ is CP10 (53)--(54), checked by Part~5.
\item \emph{Outputs used:} the tori $\varphi_\omega$, $\omega\in\Omega_*$, with (a)--(d); the measure inequality (3.1) with the \emph{absolute} constant $c_n$; $r=32\hat E\hat\rho\le C\hat\rho$ so that $r\le C\bar\gamma_2$ as in CP10 p.~561; $\abs{v-I^0}$, $\abs u$ for the reconstruction (CP10 (51), (35)).
\item \emph{Frozen frequencies and Part~4.} $H=\tilde H/2\pi$ (Lemma~\ref{bpP:p3:period-dictionary}); $\omega=\omega_0(I_0)=\omega_\mu/2\pi$ with $\omega_\mu$ the frequency map of \Pref{bpP:p4:freq}. Hence $\omega\in\mathcal D_{\gamma_1,\gamma_2,\tau}$ iff $\omega_\mu\in\mathcal D_{\alpha_1,\alpha_2,\tau}$ with $\alpha_i=2\pi\gamma_i$ (that is the dictionary of \Pref{bpP:p4:freq}), and $D\setminus D^{\rm dio}$ of (3.1) is the set Part~4 estimates; (3.1) with the \emph{same} $c_n=1+(1+2^n)(1+(4/3)^n)$ and $\meas\T^n=1$ is exactly the right side consumed by \Pref{bpP:p4:budget}, whose radius $r_K$ is $r=32\hat E\hat\rho$ here. Part~4 needs $\tau\ge n+1$; this is within the hypothesis $\tau>n$ of the proposition (take $\tau\ge n+1$ in Part~5; nothing else of Part~3 depends on $\tau$ beyond $\hat c$).
\item \emph{Conversions, with the maps named} (\Pref{bpP:p1:order}, \Pref{bpP:p1:cyl-pipeline}; Lemma~\ref{bpP:p3:period-dictionary}). The torus handed to the pipeline is the \emph{unscaled} period-one torus $K_3:=\pi\circ\varphi_\omega:\R^n\to\Ph{n}$, $\pi(J,\theta)=(\theta,J)$ (frozen order); no factor $2\pi$ is applied to it. The $2\pi$ enters exactly once, in the pipeline's chart $\Psi=\Phi\circ S_{2\pi}$ ($\Phi$ the original reconstruction map of Part~5/Part~2): $K=\Psi\circ K_3=\Phi\circ S_{2\pi}\circ\pi\circ\varphi_\omega$, with frozen frequency $w=\omega$ (CP's angular frequency is $2\pi w$), and the same $w$ is used in both the pipeline and $\mathcal D_{\gamma_1,\gamma_2,\tau}$. The map $Q$ of Lemma~\ref{bpP:p3:period-dictionary} describes the same torus in CP coordinates, $\tilde\varphi=Q\circ\varphi_\omega(\cdot/2\pi)=S_{2\pi}\circ(\cdot)$ in Part~1's language; it must not be fed to the pipeline as $K_3$ (this would scale the angles twice). The output angle $\tilde u=2\pi u$ is used only when estimates are quoted in CP's coordinates.
\end{enumerate}
\emph{Not conclusions of Part~3} (obligations of Part~5): the complex neighbourhood and uniform bounds from the literal premises (H1)--(H6), the injectivity/nondegeneracy of $\omega_\mu$ on a complex neighbourhood, the volume losses, the reconstruction of the oscillators and the seven clauses of \Pref{bpP:p0:cyltorus}, measurability and whole-torus saturation of $S$, positivity and the strict inequalities.
\sourcegap{CP10 Prop.~3 claims \emph{``a unique real-analytic embedding''}; translations $\vartheta\mapsto\vartheta+c$ preserve all stated properties, so uniqueness is neither proved nor claimed here. The Lipschitz continuity of $(\vartheta,\omega)\mapsto\varphi_\omega(\vartheta)$ passes through $\omega_0^{-1}$, which is globally Lipschitz only for bounded $D$ (locally, at scale $\rho/\bar M$, in general: separated action islands at $j^2$ with local frequency maps centred at $j$ and Hessian $1$ have inverse derivative $1$ but no global Lipschitz inverse; review E9).}

\subsection{Source gaps and suspected misprints (CP10, Prop.~3, Appendix~B)}\label{bpP:p3:errata}

Dispositions follow the independent review (\texttt{reviews/p3-review-1.md}): V = verified misprint, R = proof repair or convention change (not an erratum), T = tentative.
\begin{enumerate}[label=(E\arabic*)]
\item \textbf{V.} p.~554: ``$\hat U$ denote the $n\times n_1$ submatrix of $U$, i.e. the matrix with entries $\hat U_{ij}=U_{ij}$, for $n_1+1\le i\le n$, $1\le j\le n$'': size and index range disagree; the proof needs $n_2\times n$ (Remark~\ref{bpP:p3:rem-uhat}).
\item \textbf{V.} p.~565 (Lemma B.1): ``$M_0:=M$, $\hat M_0:=M$, $\bar M_0:=\bar M$''; $\hat M_0$ should be $\hat M$.
\item \textbf{V.} p.~566, B.1(i): ``$\hat\imath_j:I\in D_{j-1}\mapsto\omega_{j-1}\circ\omega_{j-1}^{-1}(I)\in D_j$'' is the identity; it should read $\omega_j^{-1}\circ\omega_{j-1}$ (cf.\ B.2, p.~568, $\hat\imath_+=\omega_+^{-1}\circ\omega|_D$).
\item \textbf{V, with scope.} (74) uses $\bar M_i/\bar M\le1$ without specifying compatible bounds (compatible bounds can be chosen: normalization, Remark~\ref{bpP:p3:rem-norm}); and (35), second line, carries a factor $\hat M/M$ for $\abs{v_2-I^0_2}$ that p.~572 (107) does not give (replacing it by $1$ gives what is proved). No counterexample to a sharper estimate is claimed.
\item \textbf{T (not verified).} p.~551 assigns the $1$-norm to the action, the angular sup metric, sup norms to $p,q$ and the entrywise sup norm to matrices; the proof uses operator norms and mixed conversions. The blueprint's constants are a legitimate \emph{convention change} (max norms, N-norm). No source constant is exhibited as false; the allegation that the printed constants ``cannot all hold'' is withdrawn.
\item \textbf{V (sharpening).} p.~569 (91) prints $\frac89(4s)^{\tau+1}\hat E^2$; the computation gives exponent $2(\tau+1)$. The printed inequality is true (weaker).
\item \textbf{R.} pp.~569--570: the mean value step for $\omega^{-1}$ needs the frequency segment inside the domain of the inverse; and Lipschitz bounds on non-convex $D$, $D_j$ need near/far pairs. Repaired by Lemmas~\ref{bpP:p3:ball}, \ref{bpP:p3:lip-sup}; composition domains stated (Lemma~\ref{bpP:p3:psi-lip}(i)).
\item \textbf{R (expository supplement).} Step 3 (pp.~572--573) does use (108) (``Using (108)''); the derivative bound on the open domain and analyticity on a fixed strip are not spelled out; supplied by Steps~\ref{stepP:p3:inv-dphi}, \ref{stepP:p3:lp-an} and Lemma~\ref{bpP:p3:fourier-compact}. Not a missing estimate.
\item \textbf{V.} pp.~555, 573: ``a unique real-analytic embedding'' is not true as stated (translations); global Lipschitz dependence on $\omega$ fails for unbounded disconnected $D$ (example above); bounded $D$ suffices.
\item \textbf{R.} p.~573 footnote 32: Kirszbraun (Euclidean) is replaced by max-norm McShane and the cube-covering bound $\meas F(A)\le\Lambda^d\meas A$ (\Kref{bp:p1:mcshane}, \Kref{bp:p1:lip-image}); $c_n$ is recomputed, not reproduced; not an erratum.
\item \textbf{V (dependency).} Prop.~1 is proved in CP10 only through the unpinned [4, Lemma~A.5] and a finite induction ``as in [4]''; supplied by Part~2. The pointer to Arnold [3, p.~144] in Remark~3 is historical (Appendix~B has its own induction) and is not a dependency.
\end{enumerate}
Two further verified suspected defects found by the review (for the orchestrator's ledger, not used here): CP10 p.~565 (70) and the following split print ``$k\in\Lambda$'' where $k\notin\Lambda$ is meant (with $\Lambda=\{0\}$ the zero mode would be divided by zero); CP10 p.~564 (66) prints the target $W_{v,\sigma}$ where the angle width is $s$, i.e.\ $W_{v,s}$.

\subsection{References and imported facts}\label{bpP:p3:forward}

Written parts cited: \Pref{bpP:p1:period-convert}, \Pref{bpP:p1:cyl-pipeline}, \Pref{bpP:p2:avg}, \Pref{bpP:p2:nfb}, \Pref{bpP:p2:split}, \Pref{bpP:p2:averaged-remainder}; track N: \Nref{bpN:p4:norm}, \Nref{bpN:p5:symplectic}, \Nref{bpN:p5:canonical}. Forward reference to an unwritten part: \Pref{bpP:p5:kam-application}. Track K: \Kref{bp:p1:cauchy-open}, \Kref{bp:p1:uniform-limits}, \Kref{bp:p1:mcshane}, \Kref{bp:p1:lip-image}, \Kref{bp:p1:ext-mvt}, \Kref{bp:p1:ext-banach}, \Kref{bp:p1:ext-fourier}, \Kref{bp:p1:ext-rectangle}.
The only ingredient proved elsewhere is the averaging step (Lemma~\ref{bpP:p3:avg}, proved in Part~2 from track~N); the rest is proved here from the external facts \ref{bpP:p3:ext-ift}--\ref{bpP:p3:ext-lip}. Formalization notes: the ball lemma needs a finite-dimensional path-continuation package; Lemma~\ref{bpP:p3:fourier-compact} is a bounded stand-alone task (no several-variable Vitali--Montel theorem is assumed).
